\section{Discussion and open problems}\label{sec:discusion}

\subsection{A numerical illustration}

In the six finite-size rows of Table~\ref{tab:ilustracion}, $f_1$ and $f_2$ come from integer counts obtained by exhaustive
enumeration. The values are rounded to six decimal places. In the final row, the existence of the common limit follows from
Theorem~\ref{thm:limite}, with $F(0)=1/b_0$. The decimal to its right is a numerical estimate of $1/b_0$ without a certified error
interval.

The estimate $F'(0)\approx0.516$ is also numerical. It suggests comparing $f_h(n)$ with $F((h-\tfrac12)/n)$ at finite sizes. The usefulness of
this comparison as a finite-size approximation is heuristic. The result proved in Corollary~\ref{cor:capas} concerns the sum over $K$
fixed layers as $n\to\infty$. The ratios in the table are obtained from the counts using equation~\eqref{eq:fh}.
The counts and the program that reproduces them are available in~\cite{Codigo}.

\begin{table}[ht] \centering \begin{tabular}{rccc}
\toprule
$n$ & $f_1(n)$ & $f_2(n)$ & $F(0)$ (numerical)\\
\midrule
10 & 0.127393 & 0.200891 & \\
11 & 0.111437 & 0.181943 & \\
12 & 0.123894 & 0.170782 & \\
13 & 0.122402 & 0.165452 & \\
14 & 0.109393 & 0.158389 & \\
15 & 0.114223 & 0.152947 & \\
$\infty$ & \multicolumn{2}{c}{common limit (Theorem~\ref{thm:limite})} & 0.080914\\
\bottomrule
\end{tabular} \caption{Exact values of $f_1$ and $f_2$ obtained by exhaustive enumeration and rounded to six decimal places, together
with the numerical value of the common limit. The value $0.080914$ is the second-order Richardson extrapolation of the values on the
grids with $128$ and $256$ points ($0.0809029$ and $0.0809114$). Using the grids with $64$ and $128$ points changes the extrapolation
by less than $10^{-8}$ (see~\cite[Section~9]{Perfil}). This is a numerical value, not certified by interval arithmetic.} \label{tab:ilustracion}
\end{table}

\subsection{Open questions}

\begin{enumerate}[1.] \item \emph{Rate.} It is unknown whether $f_h(n)-F(0)$, or $\log[f_h(n)/F((h-\frac12)/n)]$, is of order $1/n$. To
transfer from the reference law to the uniform law we use the relative entropy bound $d_n=o(n)$ of Theorem~\ref{thm:comparacion}, which
gives no rate, and fixed tolerances. These estimates do not justify interchanging the limits to deduce a rate. \item \emph{A growing
number of layers.} The constant term of $\log Q(n)$ depends on all layers. Since $F$ extends to a positive analytic function on
$[0,1]$~\cite[Theorem~8.1]{Perfil}, we can consider $\tau_{n,h}=n\log[f_h(n)/F((h-\frac12)/n)]$ for all layers, $1\le h\le n$. One
sufficient approach would be to prove convergence of the averages $\frac1n\sum_h\tau_{n,h}$ over those $h$ with $h/n$ in each compact
subset of $(0,1)$ and to control both endpoints uniformly. At the full-board endpoint, for example, this would require \[
 \lim_{x_0\downarrow0}\ \limsup_{n\to\infty}\ \sum_{h\le x_0n}
 \Bigl|\log\frac{f_h(n)}{F((h-\frac12)/n)}\Bigr|=0 .
\] Convergence for each fixed $h$ does not imply this condition. \item \emph{Joint law of the holes.} For $h>1$ we identify only the
mean measure of the empty crossings. Are the positions of the empty rows and columns asymptotically independent? Does the unweighted
$d_{h,1}/n$ concentrate? \item \emph{Effective constants.} The proof gives no explicit threshold $n_0$ or useful constants for boards
of moderate size. \end{enumerate}

Returning to the board with $n-1$ queens, the probability that the last queen can be placed at the unique empty crossing tends to
$1/b_0$. If $h$ queens are missing, with $h$ fixed, the same value describes the probability of adding a queen at a randomly chosen
crossing. When $h$ grows, the remaining problem is to control how the differences between the fraction $f_h(n)$ in each layer and the
corresponding value of $F$ accumulate.
